\documentclass[11pt]{article}
\usepackage{mathblatt}
\makeatletter
\setlength{\mbpfbe}{0mm}% keine Punkte (Bauregel 6.4)
\makeatother
% Rechenraum als Karo (Bauregel 6.9): n Rechenschritte = n mal 10 mm
\newcommand{\karo}[1]{\par\smallskip\noindent\begin{tikzpicture}
  \clip (0,0) rectangle ({\linewidth-1mm},{#1*10mm});
  \draw[black!16,line width=0.3pt,step=5mm] (0,0) grid ({\linewidth},{#1*10mm});
  \end{tikzpicture}\par}
\newcommand{\linie}[1]{\,{\color{black!45}\rule[-1pt]{#1}{0.4pt}}\,}
\newcommand{\pkt}{\textperiodcentered{}}
\begin{document}

\pfheftstil
\pflaufkopf[P10 \pkt{} mehr]{Seite oder Winkel mit sin, cos, tan \pkt{} Lösungen}{PG60-2}\pfseitevon
\pfheftkopf{Seite oder Winkel mit sin, cos, tan}{Lösungen \pkt{} mehr}
\begin{pfloesungb}{}
\lzb{1.}{}{AB $\approx$ 1118 m}{AB = 1400 · sin 53° $\approx$ 1118,1 m $\approx$ 1118 m}
\end{pfloesungb}
\begin{pfloesungb}{}
\lzb{2.}{}{AD $\approx$ 842,5 m}{AD = 1400 · cos 53° $\approx$ 842,5 m}
\end{pfloesungb}
\begin{pfloesungb}{}
\lzb{3.}{}{tan $\alpha$ = r/s}{tan $\alpha$ = r/s}
\end{pfloesungb}
\begin{pfloesungb}{}
\lzb{4.}{}{x = 18}{x = 9 : sin 30° = 9 : 0,5 = 18}
\end{pfloesungb}
\begin{pfloesungb}{}
\lzb{5.}{}{sin $\alpha$ = r/t}{sin $\alpha$ = r/t}
\end{pfloesungb}
\begin{pfloesungb}{}
\lzb{6.}{}{C: sin $\beta$ = b/a (dritte Option)}{sin $\beta$ = b/a (dritte Option)}
\end{pfloesungb}
\begin{pfloesungb}{}
\lzb{7.}{}{sin $\beta$ = r/t}{sin $\beta$ = r/t}
\end{pfloesungb}
\begin{pfloesungb}{}
\lzb{8.}{}{$\alpha$ $\approx$ 54,0°}{Unterschied der Seiten 14 $-$ 6 = 8 cm; tan $\alpha$ = 11 : 8 $\Rightarrow$ $\alpha$ = tan$^{-1}$($\tfrac{11}{8}$) $\approx$ 54,0°}
\end{pfloesungb}
\begin{pfloesungb}{}
\lzb{9.}{}{$\alpha$ $\approx$ 65°}{cos $\alpha$ = 5 : 12 $\Rightarrow$ $\alpha$ = cos$^{-1}$($\tfrac{5}{12}$) $\approx$ 65,4° $\approx$ 65°}
\end{pfloesungb}
\begin{pfloesungb}{}
\lzb{10.}{}{Berghöhe $\approx$ 61,0 m; D liegt auf 61,0 + 10,0 $\approx$ 71,0 m}{BC = 118,8 · sin 30° $\approx$ 59,4 m; Berghöhe = BC + 1,6 $\approx$ 61,0 m; D liegt auf 61,0 + 10,0 $\approx$ 71,0 m}
\end{pfloesungb}
\begin{pfloesungb}{}
\lzb{11.}{}{PB $\approx$ 1,85 m}{PB = 1,35 : tan 36,1° $\approx$ 1,85 m}
\end{pfloesungb}
\begin{pfloesungb}{}
\lzb{12.}{}{BD $\approx$ 440 m}{$\angle$CAD = 90° $-$ 52° = 38°; $\angle$DAB = 71° $-$ 38° = 33°; BD = 677 · tan 33° $\approx$ 440 m}
\end{pfloesungb}
\end{document}